\documentstyle[12pt,fullpage,steve]{article}
\setcounter{page}{1}

\begin{document}

\title{EE 150 Quiz \#1}
\author{Due: Oct. 7 (NOW), 1996 2:45 pm (in class)}
\date{}
\maketitle

As announced, this is the first of a series of in-class quizes. Some
will be done in-class, while others will be due over a 1 week period.
Grades will be give for these quizes, and they will count towards your
final grade.  These are in addition to ongoing assignments. 

\section*{Question 1 of 1}

Tell us what the following program fragment does.  To determine
the answer, work through the code by hand and show the output that
results.  {\bf You must show how these programs execute with the
``tables'' model shown in class  in order to get credit.} Hint:  A
table method of showing the output involves showing the time or
execution steps on one axis (say the horizontal axis) and each
variable that changes on the other axis (say the vertical axis).  Hint
2: It wouldn't hurt, after the table, to explain your reasoning if you
feel your table is 
not clear.

\begin{verbatim}
main()
{
  int a;  int b;   int c;  int d;  double dd;

  int    dostuff(int,int,int);
  double dostuff2(int,int,int);

  a = 2;  b = a * 3 - 4;  c = b - a + 5 * 1.2;  d = dostuff(c, b, a);
  printf("xx d = %i\n", d);  printf("xx b = %i\n", b);
  a = b - d;
  printf("xx a = %i\n", a);  printf("xx b = %i\n", b);
  dd = dostuff2(c,b,a);      printf("xx dd : %5.2f\n", dd);
  dd = (int) dd / 2;         printf("xx dd : %5.1f\n", dd);
}

int dostuff(int c, int a, int b)
{
  printf("y0: a b c = %i,%i,%i\n", a, b, c);
  return a + b * c;
}

double dostuff2(int c, int a, int b)
{
  a = a * b;
  printf("y1: a b c = %5.1f,%i,%i\n", (double) a, b, c);
  return (double) a + c * b;
}
\end{verbatim}

\end{document}
